\documentclass[10pt,twocolumn]{article}

% ------------------ PAQUETES ------------------
\usepackage[spanish]{babel}
\usepackage[utf8]{inputenc}
\usepackage[T1]{fontenc}
\usepackage{geometry}
\usepackage{graphicx}
\usepackage{amsmath, amssymb}
\usepackage{siunitx}
\usepackage{authblk}
\usepackage{abstract}
\usepackage{caption}
\usepackage{titlesec}
\usepackage{multicol}
\usepackage{hyperref}
\usepackage{currfile}
\usepackage{float}

\geometry{margin=1.8cm}

% ------------------ FORMATO DE SECCIONES ------------------
\titleformat{\section}{\large\bfseries}{\thesection}{1em}{}
\titleformat{\subsection}{\normalsize\bfseries}{\thesubsection}{1em}{}

% ------------------ TÍTULO ------------------
\title{\textbf{Aproximación manual a la fotometría de apertura}\\Trabajo sobre el cúmulo abierto Messier 41\\Tercera parte: Introducción al machine learning}


% ------------------ AUTORES ------------------
\author[1]{Bruno Abello Laurent}

\affil[1]{Universidad de Los Andes, Departamento de Física, Colombia}

% Correos (puedes modificar)
\renewcommand\Authands{, }
\date{}

\begin{document}

\twocolumn[
\maketitle
\begin{center}
\small
\textit{Correo institucional:} \href{mailto:{b.abello@uniandes.edu.co}}{b.abello@uniandes.edu.co} \\
\textit{Repositorio de GitHub: } \href{https://github.com/Brunelleschk1/Personal-projects/tree/main/Proyectos\%20de\%20Astronom\%C3\%ADa/Laboratorio\%203}{Repositorio del laboratorio}
\end{center}

\begin{abstract}
\noindent

\end{abstract}

\noindent\textbf{Palabras clave:} Messier 41, HDBSCAN, Vector Point Diagram, Color Magnitude Diagram. 
\vspace{0.5cm}
]

% ------------------ CONTENIDO ------------------

\section{Introducción}

El presente informe se puede ver como la culminación de varias semanas de trabajo sobre los métodos computacionales
que se utilizaron históricamente en el estudio de los cúmulos abiertos. A través de métodos cada vez más refinados, 
nos hemos acercado al uso del machine learning para identificar cúmulos abiertos entre los datos que obtienen 
las herramientas más poderosas de la humanidad para detectar y catalogar cuerpos celestes en el universo. 

Como una prueba de lo poderoso pero poco comprendido que puede ser un método de machine learning, trabajamos
con todos los cuerpos a 100 pársecs alrededor de nuestro sol, con el fin de encontrar clusters en este espacio.
Lo interesante de este estudio es que la humanidad ya ha estudiado bien esta región, y ha encontrado dos 
cúmulos, Coma y Hyades. Nosotros no queremos limitarnos a estos dos, y buscaremos cualquier indicio de que 
un tercero o más cúmulos se encuentran en esta región. Es decir, por primera vez nos interesamos a un 
trabajo de exploración y no a uno de encontrar algo que sabemos que está ahí. 

Las explicaciones del método que se utiliza, HDBSCAN, ya se dieron en semanas anteriores, por lo que, 
para esta ocasión, podemos ser más directos con nuestros resultados y otros detalles. 
\subsection{Objetivos}

Los objetivos que tenemos, en términos simples, son los siguientes:
\begin{itemize}
    \item Implementar una manera de analizar cientos de miles de datos de GAIA
    \item Implementar formas, simples o avanzadas, de determinar qué es un cluster real
    \item Entender las limitaciones que tiene el Machine Learning 
    \item Realizar un análisis manual que de un paso más allá de lo que hace un computador.
\end{itemize}

\section{Procedimiento}

Teniendo en cuenta estos objetivos, proceder fue bastante directo: al hablar con la clase, se determinó 
que se podía cortar el cielo en 64 porciones iguales y analizar cada una por separado. Resulta que
HDBSCAN tiene un crecimiento altísimo en tiempo con la cantidad de datos que se le proporcionan, por lo que
analizar pequeñas celdas de 5000-10000 estrellas es mucho mejor que analizar todo el cielo al tiempo. 

Acto seguido, se empezaron a implementar cortes y restricciones a la función de análisis, de manera que 
no se tuviera que graficar todo el cielo por un falso cúmulo de 30000 estrellas. Además, se realizó a posteriori
una gráfica de todo el cielo, en ascención recta contra declinación, para ver en dónde y qué forma tienen los 
cúmulos detectados. 

El como se llegó a la función final y el cómo se resolvió hacer todo esto pertenece a la sección siguiente.

\section{Problemas presentados y soluciones aplicadas}

\begin{enumerate}
  \item  Uno de los primeros problemas encontrados fue qué presentar como resultado. 
  Se resolvió que lo mejor era presentar dos CMDs y dos VPDs, unos sin silueta y otros dos con ella, 
  además de un histograma con la distribución de miembros por cluster encontrado sobre la distancia. 
  \item El segundo problema fue darse cuenta de que hay muchos cúmulos encontrados por HDBSCAN que son, 
  de alguna manera, basura. Tienen 60, o 20000 miembros, y están muy lejos de ser un cúmulo real. 
  La solución, sencilla, es cortar cúmulos con menos de 90 miembros o más de unos 1050 miembros, 
  para conservar los que, al menos, tienen un semblante real. 
  \item Posteriormente, se encuentra el siguiente problema, y es que muchos cúmulos encontrados
  no se comportan, en su extensión, como un cúmulo real: son enormes. Por ende, se elimina un cúmulo
  que tenga más de 34 pársecs de extensión. 
  \item Además de los problemas anteriores, un problema más elusivo se puede encontrar en la forma de la distribución 
  del cúmulo. Se espera que sea gaussiano, pero la mayoría del tiempo se obtiene una pendiente creciente, que asumimos
  infinita en el espacio más allá de los 100 clusters: un cono de estrellas hacia el infinito. Para eliminar estos cúmulos, 
  se implementó un corte un poco curioso: calcular la pendiente de los datos y, si supera un threshold, eliminar el cúmulo. 
  \item El lector podrá comprender que esto no elimina todos los cúmulos como se quiere, y se siguen obteniendo datos. 
  Sin embargo, resultó que, para el estudiante, implementar cortes automáticos correctos empezó a resultar complicado
  y poco fructífero. Por ende, decidió que iba a realizar un ajuste gaussiano a cada cúmulo restante, para luego revisar 
  manualmente qué tan convincente era el ajuste. 
  \item Por aparte, también hay que notar algo muy importante, y es que los datos que se utilizaron no incluían las regiones
  de Coma y Hyades. Por ende, estos datos se volvieron a añadir gracias a un código de python, y se sumaron al análisis. 
  De esta manera, se sabe cómo debería lucir un cúmulo real con la presentación que se instauró para hacer el análisis. 
\end{enumerate}

Nótese, primeramente, que el volver a añadir los datos de Hyades y Coma resultó en la superposición de regiones en el mapa.
Por ende, se ven en el mapa dos regiones especiales, lo que no es muy grave pero si amerita una explicación previa al análisis
serio de los datos. 
Por otro lado, y como dato o consejo para el lector, se ha de entender que los pasos propuestos en la lista anterior están aplicados 
en el código en ese orden exactamente. Esto es debido a que los cortes fuertes eliminan lo realmente innecesario, y lo difícil de procesar, 
como un cúmulo muy extenso o un cúmulo con 20000 miembros, dejando así los métodos más exigentes, como el cálculo de pendiente o el 
ajuste de gaussiana, a una serie reducida, y más real, de cúmulos. Así, el código es mucho más rápido. 

\section{Resultados}

Como podrá entender el lector, los resultados de este informe son, en cierta medida, subjetivos. Sí, hay muchos criterios 
que nos permiten decidir si algo es o no es un cúmulo de manera objetiva, pero lo cierto es que no todo se resuelve con código. 
Esta sección es, entonces, una presentación sucinta, comparativamente hablando, de los resultados encontrados. 
O sea, de lo más prometedor. 

Primeramente, observemos los resultados de los clusters conocidos: Coma y Hyades. Los resultados obtenidos se presentan en 
PONER Y REFERENCIAR GRÁFICAS DE COMA Y HYADES. Son, adicionalmente, las regiones especiales del cielo en la figura PONER
EL CIELO COMPLETO. Se observa de los CMDs, VPDs y distribuciones de distancia el comportamiendo estándar de un cúmulo bajo
las condiciones impuestas por el código. Los comportamientos de la distribución del cúmulo en función de la distancia son 
normales, se observa muy poco ruido en los CMDs, y una región bien comportada en los VPDs. Curiosamente, en el estudio de 
EL QUE ESTÁ MÁS LEJOS, se obtiene un cúmulo más pequeño. Sin embargo, por el número de miembros encontrado, CITAR CUANTOS SON es probable que este
cumulito haga parte del cúmulo más grande ya conocido.

El trabajo más importante, sin embargo, es preguntarnos si existen cúmulos en esta esfera alrededor del sol que no conocemos.
Ya se explicaron los cortes automáticos a los cúmulos encontrados por machine learning, y ahora toca explicar aquellos que 
se juzgan manualmente. Nótese que de 64 regiones, las regiones 2, 7, 14, 19, 20, 22, 25, 31, 39, 42, 44, 47, 48, 51, 52 y 62 
son las únicas que presentan algún tipo de resultado medianamente coherente. O sea 16 regiones de 64. De estas 16 regiones, visibles
en la figura CITAR MAPA, varias pueden ser descartadas por una u otra razón. 

Hay que empezar estos descartes por notar que a excepción de Coma y Hyades, ninguno de los cúmulos encontrados tienen una distribución en 
RA y DEC similar a la de un cúmulo. Todas las estrellas están dispersas a lo largo de la celda, por lo que primero hablaremos de las anomalías
de este estilo: Las regiones 19, 20 y 44. La concentración en 19 está acumulada hacia la celda 20, por lo que juntos podrían ser un semblante de cúmulo.
Sin embargo, no se encuentran en el mismo punto del espacio, ya que uno está centrado hacia los 65 pársecs y otro hacia los 85. Además, el cúmulo en la celda
19 es más extenso que 30 pársecs, y su CMD tiene mucho ruido en la zona inferior, junto con una secuencia principal pobre. Además, la distribución en el mapa estelar
de la celda 20 no tiene coherencia con la de la celda 19. En definitiva, juntas no forman un cúmulo, y por separado 19 tiene demasiado ruido y 20 tiene una mala distribución
en distancia, sin hablar de que la mayoría parece ser ruido de GAIA, puesto que si fuera un cúmulo joven, no estaría tan esparcido en el mapa estelar. 
PONER IMAGENES DE 19 Y 20, JUNTO CON LA CITA DEL MAPA ESTELAR. TAL VEZ PONER DISTRIBUCIONES RADEC DE COMA Y HYDES SEA BUENA IDEA???

La celda 44 muestra una acumulación de estrellas enorme cerca a las celdas 41, 53 y 52. Sin embargo, estas celdas no tienen cúmulos coherentes con los de la celda 44, y mientras que 
el cúmulo 3 encontrado en esta celda es enorme, con una mala distribución en distancia y muchísimo ruido en el CMD, el primer cúmulo no es muy diferente: es casi demasiado pequeño, 
pero está muy bien distribuido en distancia y en el VPD. Además, PONER L DISTRIB DE 44_1 EN FULL, la distribución en el mapa espacial es demasiado extensa para un cúmulo recién nacido.
Es decir, la celda 44 sólo contiene ruido y ni la celda 44, ni la 19 ni la 20 encontraron algún cúmulo prometedor. Con eso, podemos movernos a las siguientes 13 regiones, 
todas mal distribuidas en el mapa estelar; se van a discutir de menos prometedores a más prometedores, no por orden ascendente o descendente.

Empezemos esta siguiente parte por la celda número 7. Ambos tienen ruido en el CMD, pero una buena secuencia principal.




\begin{figure}[H]
    \centering
    \includegraphics[width=0.8\linewidth]{../Resultados_completos.png}
    \caption{Compilación de los CMDs obtenidos con el estudio de los datos mediante HDBSCAN. El cluster Messier 41 es,
    se espera, el cluster 1 tal y como lo marca el método.}
    \label{resultadosCompletos}
\end{figure}

\section{Conclusiones}

\section*{Comentarios al margen}

\bibliographystyle{apalike}
\bibliography{refs}
\end{document}